% 哥白尼原则（综述）
% license CCBYSA3
% type Wiki

本文根据 CC-BY-SA 协议转载翻译自维基百科\href{https://en.wikipedia.org/wiki/Copernican_principle}{相关文章}。

在物理宇宙学中哥白尼原理（Copernican principle）指出：人类并非宇宙中的特权观察者，\(^\text{[1]}\)来自地球的观测结果可视为来自宇宙中平均位置的代表性观测。该原理以哥白尼的日心说（heliocentrism）命名，是一种工作假设，其思想源自对哥白尼“地球运动”论证的宇宙学延伸。\(^\text{[2]}\)
\subsection{起源与意义}
赫尔曼·邦迪（Hermann Bondi）在二十世纪中期将该原理命名为“哥白尼原理”，但该思想本身可追溯至十六至十七世纪期间从托勒密体系（Ptolemaic system）向新范式的转变——后者不再将地球置于宇宙中心。哥白尼提出，行星的运动可通过假设太阳处于静止的中心位置来解释，从而与地心说形成对比。他认为，行星的视逆行运动是一种由地球绕太阳运动所造成的幻象，而在哥白尼模型中，太阳被置于宇宙的中心。值得注意的是，哥白尼的主要动机是出于对旧体系的技术性不满，而非出于支持任何“平庸原理”（principle of mediocrity）的哲学立场。\(^\text{[3]}\)

尽管哥白尼的日心模型常被描述为“将地球降级”——使其失去在托勒密地心模型中的中心地位，但真正采纳这一新宇宙观的，是哥白尼的继承者，尤其是十六世纪的焦尔达诺·布鲁诺（Giordano Bruno）。在传统观念中，地球的中心位置被认为是“最低且最污秽之处”；而伽利略则指出，地球实际上参与了“群星的舞蹈”，而非停留在“宇宙污秽与浮渣的积聚之所”。\cite{ref4, ref5}  

至二十世纪后期，卡尔·萨根（Carl Sagan）以诗意的语言进一步强化了哥白尼原理的精神，他问道：“我们是谁？我们发现自己居住在一颗微不足道的行星上，这颗行星环绕着一颗平凡的恒星，处于银河系的某个被遗忘的角落，而宇宙中星系的数量远远超过人类的数量。”\cite{ref6}
